\chapter{Background and Related Works} \label{ch:state-of-the-art}

\section{Technical Foundations}\label{sec:tech_foundations}

\subsection{Robotic insertion as a contact-rich manipulation problem}

Robotic insertion is a common assembly task in which one part is moved into a corresponding opening, such as a peg into a hole or a component into a fixture. Although the final state is geometrically defined, the execution is not purely a positioning problem. Small deviations in position or orientation can cause contact between the parts \cite{qiaoRoboticPegHoleInsertion1993,xuCompareContactModelbased2019}, and the resulting interaction forces can determine whether the insertion succeeds or fails \cite{raibertHybridPositionForce1981,zhaoPeginHoleAssemblyBased2020}.

For this reason, robotic insertion can be understood as a contact-rich manipulation task. During the insertion motion, the robot may transition from free-space movement to one-sided or multi-point contact. Larger pose misalignments can lead to multi-point contact conditions \cite{tangAutonomousAlignmentPeg2016}, which may require explicit alignment or compliant control strategies before a reliable insertion can be achieved \cite{raibertHybridPositionForce1981,xuCompareContactModelbased2019}. Depending on geometry, tolerance, surface properties, and alignment, contact can guide the part into the opening or lead to sliding, scraping, wedging, or jamming \cite{tangAutonomousAlignmentPeg2016,feiContactJammingAnalysis2004b,beltran-hernandezLearningForceControl2020}.

This makes insertion relevant for sensor-based classification. Robot pose alone may not capture the physical interaction, while force and torque signals provide information about the contact state and the progression of the insertion over time \cite{qiaoRoboticPegHoleInsertion1993,jasimContactstateMonitoringForceguided2014,magistrisExperimentalForceTorqueDataset2018}.

In this thesis, insertion is therefore treated as a practical contact-rich process that produces characteristic multivariate time-series signals during successful and failed attempts \cite{magistrisExperimentalForceTorqueDataset2018,beltran-hernandezLearningForceControl2020}.

\subsection{Sensor information in robotic insertion}

Because the physical interaction during insertion changes over time, the task naturally produces sequential sensor information. Relevant signals can include force and torque measurements, robot poses, velocities, joint states, gripper states, or external measurements. Six-axis force/torque sensors measure forces and moments along multiple axes, which makes them suitable for observing mechanical interaction between the inserted part and the environment \cite{liMultiaxisForceTorque2025,qiaoRoboticPegHoleInsertion1993,tangAutonomousAlignmentPeg2016}. In insertion-specific learning datasets, force/torque and pose signals have therefore been recorded as central data sources \cite{magistrisExperimentalForceTorqueDataset2018}.

The temporal development of these signals can contain information that is not visible in a single time step. Contact onset, sliding, correction motions, jamming, or successful seating can appear as characteristic changes in the recorded trajectories. Related work uses wrench and pose data for contact-state monitoring \cite{jasimContactstateMonitoringForceguided2014} and contact-pose identification in peg-in-hole assembly \cite{jinContactPoseIdentification2021a}, which supports the relevance of sequential sensor information for interpreting the insertion process.

Robotic manipulation datasets increasingly follow this sequential and multimodal view. Insertion-specific work has recorded force/torque and pose data for learning from multi-shape insertion attempts \cite{magistrisExperimentalForceTorqueDataset2018}. Broader contact-rich manipulation and assembly datasets combine modalities such as vision, proprioception, force, audio, or action information \cite{sliwowskiREASSEMBLEMultimodalDataset2025,fangRH20TComprehensiveRobotic2024}. Recent work on auditory or audio-visual learning further shows that contact-rich manipulation can benefit from modalities beyond vision and robot state alone \cite{thankarajThatSoundsRight2023,mejiaHearingTouchAudioVisual2024b}.

For this thesis, each insertion attempt is therefore represented as a multivariate time series. Each sensor channel describes one aspect of the robot or process state, while the sequence describes how these quantities evolve during the attempt. This representation preserves both the multi-channel structure and the temporal information that may distinguish successful from failed insertions.

\subsection{Formulation as a multivariate time-series classification problem}

Time-series classification describes the task of assigning a discrete label to an ordered sequence of measurements \cite{fawazDeepLearningTime2019}. In the multivariate case, each sample contains several synchronized channels rather than a single signal \cite{bagnallGreatTimeSeries2017,ruizGreatMultivariateTime2021}. A recorded insertion attempt can therefore be written as a sequence $X \in \mathbb{R}^{T \times C}$, where $T$ denotes the number of time steps and $C$ the number of sensor channels. The corresponding class label $y$ describes the outcome of the attempt, for example successful or failed insertion.

This formulation is suitable for robotic insertion because the relevant information is often distributed over the complete signal trajectory. A classifier does not only evaluate individual force, torque, or motion values, but learns patterns in their temporal evolution and in their combination across channels. This distinguishes the task from a static threshold-based decision, where only selected values or manually defined features are considered \cite{fawazDeepLearningTime2019}. The same principle is reflected in robotic assembly work that uses temporal wrench and pose information for contact-state or contact-pose recognition \cite{jasimContactstateMonitoringForceguided2014,jinContactPoseIdentification2021a}.

The classification perspective also provides a structured way to evaluate generalization. Standard time-series classification archives and benchmark studies emphasize the importance of clearly defined train/test splits and comparable evaluation procedures \cite{dauUCRTimeSeries2019,bagnallGreatTimeSeries2017,ruizGreatMultivariateTime2021}. In robotic insertion, trials can additionally be grouped by experimental conditions such as geometry, tolerance, gripper configuration, or imposed deviations. Holding out such groups makes it possible to test whether a model has learned general sensor patterns of successful and failed insertions, or whether its performance depends strongly on conditions already seen during training.

In this thesis, the insertion data are therefore treated as a supervised multivariate time-series classification problem. Each trial is represented by synchronized sensor channels over a fixed observation window, and the classifier is trained to predict the binary insertion outcome. This formulation connects the physical insertion process to established time-series classification methods and provides the basis for the benchmark design used in the following chapters.

\section{Related work on time-series classification}

Time-series classification has been studied with a wide range of methods, from distance-based classifiers to feature-based approaches, ensemble methods, convolutional neural networks, and transformer-based models. Earlier benchmark studies showed that no single representation is optimal for all datasets, which motivated methods that combine several views of the same time series \cite{bagnallGreatTimeSeries2017,dauUCRTimeSeries2019}. More recent work has also shown that deep learning and randomized convolutional transformations can achieve highly competitive performance while reducing the need for manually designed features \cite{fawazDeepLearningTime2019}, with examples being ROCKET \cite{dempsterROCKETExceptionallyFast2020} and InceptionTime \cite{ismailfawazInceptionTimeFindingAlexNet2020}.

\subsection{ROCKET \cite{dempsterROCKETExceptionallyFast2020}}

ROCKET is based on a different idea. Instead of learning convolutional filters through gradient-based training, it applies a large number of randomly generated convolutional kernels to the input time series. From the transformed signals, simple features such as the maximum value and the proportion of positive values are extracted. These features are then passed to a linear classifier, typically ridge regression or logistic regression.

The relevance of ROCKET lies in the combination of simplicity, speed, and strong classification performance. Dempster et al. show that random convolutional kernels can achieve highly competitive results with substantially lower computational cost than many previous state-of-the-art methods. This makes ROCKET an important baseline for applied time-series classification problems. For robotic insertion data, ROCKET is relevant because force, torque, and motion signals may contain local temporal patterns that can be captured by convolutional filters, while the linear classifier keeps the final model comparatively simple.

\subsection{InceptionTime \cite{ismailfawazInceptionTimeFindingAlexNet2020}}

InceptionTime represents a deep convolutional approach to time-series classification. It adapts the Inception architecture, originally developed for image recognition, to one-dimensional temporal signals. The model uses several convolutional filters with different kernel sizes in parallel, allowing it to extract temporal patterns at multiple scales. Several Inception modules are stacked and combined in an ensemble to improve robustness and accuracy.

This architecture is relevant because robotic insertion signals may contain both short events and longer process-level patterns. Short contact events, oscillations, or force peaks may occur within a small time window, while the overall progression from approach to contact and final insertion unfolds over a longer sequence. A multi-scale convolutional model such as InceptionTime is therefore a meaningful deep learning baseline for this type of data. InceptionTime also provides a useful contrast to transformer-based models, because it learns local and multi-scale temporal features through convolutions rather than attention mechanisms.

\subsection{HIVE-COTE 2.0 \cite{middlehurstHIVECOTE20New2021}}

HIVE-COTE 2.0 is a heterogeneous ensemble for time-series classification. Instead of relying on one representation of the signal, it combines classifiers that operate on different transformed views of the data. These include shapelet-based, dictionary-based, interval-based, and convolution-based components. This design reflects the observation that discriminative information in time series can appear in different forms, such as local subsequences, repeated symbolic patterns, temporal intervals, or phase-dependent shapes.

The main strength of HIVE-COTE 2.0 is its classification accuracy. Middlehurst et al. report that HIVE-COTE 2.0 achieved state-of-the-art average performance on both the univariate UCR archive and the multivariate UEA archive. For this reason, it is an important reference point when discussing the upper end of classical time-series classification performance. However, this accuracy comes at the cost of considerable computational complexity and a relatively large ensemble structure. In the context of this thesis, HIVE-COTE 2.0 is therefore most useful as a high-performance literature reference, while simpler and more scalable methods may be more practical for repeated benchmarking on robotic insertion data.

\subsection{PatchTST \cite{nieTimeSeriesWorth2023b} and transformer-based time-series models}

Transformer models were originally introduced for sequence modeling with attention mechanisms \cite{vaswaniAttentionAllYou2017}. In time-series learning, attention-based models are attractive because they can, in principle, relate information from distant time steps without relying on recurrent processing. However, applying standard transformers directly to long time series can be computationally expensive, since the attention mechanism scales with the sequence length.

PatchTST addresses this issue by dividing a time series into patches before applying the transformer encoder. Instead of treating every individual time step as one token, short subsequences are used as input tokens. This reduces the effective sequence length and allows the model to process longer temporal contexts more efficiently. PatchTST also uses a channel-independent design, where each variable is processed as a separate univariate sequence with shared model weights. This design is particularly relevant for multivariate sensor data, because it separates temporal representation learning within each channel from later aggregation across channels.

The original PatchTST work focuses primarily on long-term forecasting and self-supervised representation learning rather than classical time-series classification. Nevertheless, the architectural ideas are relevant for classification tasks with long multivariate sequences. In the context of robotic insertion, patching can summarize local signal segments such as contact transients or short force responses, while the transformer encoder can model how these segments evolve over the full insertion attempt. For this reason, PatchTST is used in this thesis as the main model family for classifying successful and failed insertion attempts from multivariate sensor time series.

\subsection{Relevance for this thesis}

The reviewed methods represent different perspectives on time-series classification. HIVE-COTE 2.0 provides a strong accuracy-oriented ensemble reference \cite{middlehurstHIVECOTE20New2021}. ROCKET offers an efficient randomized convolutional baseline \cite{dempsterROCKETExceptionallyFast2020}. InceptionTime represents a deep convolutional approach for multi-scale temporal feature extraction \cite{ismailfawazInceptionTimeFindingAlexNet2020}. PatchTST introduces transformer-based sequence modeling with patching and channel handling that are well aligned with long multivariate sensor recordings \cite{nieTimeSeriesWorth2023b}.

This thesis focuses on PatchTST because the insertion recordings are long, multichannel sequences in which relevant information may be distributed across the entire attempt. At the same time, the other methods provide important context for interpreting the results. They show that strong time-series classification performance can be obtained through several different model families, and they motivate the use of baselines when evaluating whether a transformer-based model is necessary for the proposed robotic insertion dataset.

\section{Related robotic insertion and assembly datasets}

The dataset contribution of this thesis is best understood in relation to existing data resources for robotic insertion, contact-state recognition, and contact-rich assembly. These works differ substantially in task scope, sensing modalities, label structure, and intended use. Some datasets focus directly on peg-in-hole insertion and force/torque sensing, while others address broader assembly or manipulation tasks with multimodal recordings. For this thesis, the most relevant comparison is therefore not only the size of a dataset, but also whether it supports outcome classification, controlled variation of insertion conditions, and benchmark-style evaluation across different task distributions.

\subsection{Insertion-specific force and contact datasets}

The closest insertion-specific reference is the experimental force-torque dataset for robot learning of multi-shape insertion by De Magistris et al. The dataset contains peg-in-hole insertion data with force/torque and pose information for several convex peg shapes \cite{magistrisExperimentalForceTorqueDataset2018}. Its purpose is to support learning from physical interaction data, including learning insertion behavior from six-axis force/torque measurements and related contact tasks such as shape recognition \cite{magistrisExperimentalForceTorqueDataset2018}. This makes it highly relevant to the present thesis, because both works treat force-based insertion as a data-driven learning problem. However, the focus differs. De Magistris et al. emphasize learning insertion behavior and shape-related contact information, whereas this thesis focuses on classification of successful and failed insertion attempts under controlled variations of geometry, tolerance, gripper configuration, and imposed deviations.

Related work on contact-state monitoring also records data that are conceptually close to insertion datasets, even when the main contribution is a recognition method rather than a reusable public benchmark. Jasim and Plapper capture wrench and pose signals for different contact states or phases in force-guided peg-in-hole assembly \cite{jasimContactstateMonitoringForceguided2014}. These signals are then used to build models for recognizing the current contact state. This supports the idea that force/torque and pose trajectories contain discriminative information about the assembly process. At the same time, the task formulation is different from the one used in this thesis. Contact-state monitoring aims to recognize intermediate physical states during the assembly, while this thesis uses complete insertion attempts as samples and assigns an outcome label to each trial.

A similar distinction applies to contact-pose identification. Jin et al. classify contact poses in peg-in-hole assembly from a sequence of contact measurements \cite{jinContactPoseIdentification2021a}. Their work is relevant because it explicitly uses temporal contact information to infer the relative pose between peg and hole under uncertainty. However, the classification target is the contact pose used for alignment and control, not the final success or failure of an insertion attempt. In contrast, the dataset in this thesis is designed around complete trials and evaluates whether time-series models can distinguish successful from failed insertions, including cases from unseen or shifted experimental conditions.

\subsection{Broader contact-rich assembly datasets}

Beyond insertion-specific work, several recent datasets address contact-rich manipulation and assembly on a broader scale. REASSEMBLE is a multimodal dataset for robotic assembly and disassembly built around the NIST Assembly Task Board 1 \cite{sliwowskiREASSEMBLEMultimodalDataset2025}. It includes actions such as pick, insert, remove, and place, and records multimodal sensor data including multi-view RGB cameras, event cameras, force/torque sensing, microphones, and robot proprioception \cite{sliwowskiREASSEMBLEMultimodalDataset2025}. The dataset also provides annotations for tasks such as temporal action segmentation, motion policy learning, and success or anomaly detection \cite{sliwowskiREASSEMBLEMultimodalDataset2025}. REASSEMBLE is therefore very relevant as a modern contact-rich assembly dataset, but its scope is broader than the focused insertion setting of this thesis. It captures complete assembly and disassembly behavior across several objects and actions, while this thesis isolates insertion attempts and studies their classification under controlled experimental variation.

RH20T represents an even broader view of robot manipulation data. It contains a large number of real-world contact-rich manipulation sequences across diverse tasks, robots, and viewpoints, with modalities such as vision, force, audio, action information, and human demonstration videos \cite{fangRH20TComprehensiveRobotic2024}. Such a dataset is important for general robot learning and multimodal representation learning. However, its purpose is not to provide a targeted benchmark for insertion outcome classification. Compared with RH20T, the dataset in this thesis is much narrower in task scope, but more specific in the controlled variation of insertion-relevant factors and in the evaluation of success/failure classification.

FurnitureBench is another relevant assembly benchmark, but with a different emphasis. It focuses on real-world long-horizon furniture assembly and provides pre-collected demonstrations, reproducible physical furniture models, and a simulation environment \cite{heoFurnitureBenchReproducibleRealworld2025a}. Its strength lies in benchmarking complex manipulation policies over extended assembly sequences. The relevance to this thesis is that furniture assembly also contains contact-rich subtasks such as insertion, alignment, and fastening. Nevertheless, FurnitureBench is not primarily designed as a sensor-level time-series classification dataset for individual insertion attempts. It therefore complements, rather than replaces, a focused insertion dataset.

\subsection{Comparison dimensions for this thesis}

The reviewed datasets show that robotic assembly data can be collected for different purposes. Some works focus on contact-state or contact-pose recognition during insertion \cite{jasimContactstateMonitoringForceguided2014,jinContactPoseIdentification2021a}. Others provide broader multimodal datasets for manipulation, assembly, or policy learning \cite{sliwowskiREASSEMBLEMultimodalDataset2025,fangRH20TComprehensiveRobotic2024,heoFurnitureBenchReproducibleRealworld2025a}. The dataset developed in this thesis is positioned between these directions. It is narrower than large-scale general manipulation datasets, but more directly structured around the question of whether insertion outcomes can be classified from multivariate sensor time series.

For this reason, the most important comparison dimensions are task specificity, sensing modalities, label type, controlled variation, and evaluation protocol. Task specificity describes whether the data focus on insertion itself or on broader manipulation and assembly. Sensing modalities describe whether the dataset contains force/torque, robot state, vision, audio, or other signals. Label type distinguishes between contact states, contact poses, action phases, demonstrations, anomalies, and final trial outcomes. Controlled variation is especially important for this thesis, because geometry, tolerance, gripper configuration, and positional or rotational deviations can be used to define in-distribution and out-of-distribution evaluation settings. Finally, the evaluation protocol determines whether a dataset supports only random trial splits or also more demanding group-based tests in which entire experimental conditions are held out.

From this perspective, the contribution of this thesis is not only the collection of another set of insertion trials. The intended contribution is a dataset and benchmark structure that connects contact-rich robotic insertion with supervised multivariate time-series classification. By recording complete insertion attempts with synchronized sensor channels and by organizing the data according to controlled experimental factors, the dataset enables both classification of insertion outcome and analysis of generalization across task variations. This focus distinguishes it from broader manipulation datasets and from insertion studies that primarily target control, alignment, or intermediate contact-state recognition.